\section{Introduction}

Database systems have long stopped treating multidimensional data as a
specialty. Geometric and geographic objects, ranges, full-text signatures and
vector representations are handled by general-purpose systems alongside scalar
types. No single tree structure serves all of them, and systems have responded
with frameworks rather than structures: the generalized search tree
\cite{hellerstein1995generalized} implements page layout, descent, split, the
concurrency protocol \cite{kornacker1997concurrency} and write-ahead logging,
and delegates the data-specific decisions to an \emph{operator class}.
Substituting bounding rectangles yields an R-tree \cite{guttman1984r};
substituting signatures yields a full-text index; substituting ranges yields a
range index. This is what makes an index framework extensible by users rather
than only by the authors of the system
\cite{chilingarian2004postgresql,rogov2026}.

Extensibility is bought with ignorance. The framework knows containment
(\code{consistent}), key combination (\code{union}), the cost of extending a
key (\code{penalty}) and node partitioning (\code{picksplit}), and nothing
else~--- in particular, no order on keys, and no way to derive one. Bulk
loading, which for a B-tree means sorting the input and packing pages, is
therefore unavailable, and indexes on such frameworks are built the only way
that requires no order: by inserting entries one at a time. Every insertion
descends to a leaf, so construction costs $O(N \log_f N)$ page accesses; the
accesses are random, because insertion order is unrelated to physical
placement; and pages end up filled to about 70\,\%, because they are produced
by splits rather than by packing.

The gap is conspicuous because bulk loading of R-trees is well understood.
Packed R-trees \cite{roussopoulos1985direct}, Hilbert packing
\cite{kamel1993hilbert} and sort-tile-recursive packing
\cite{leutenegger1997str} all sort by some function of the key and fill nodes
consecutively. None of them applies here: each is formulated in terms of
coordinates or of the area of a bounding rectangle, and an operator class over
signatures or over ranges offers neither.

\paragraph{Problem.} We therefore ask the question one level up. Given an
extensible index framework, (i)~what is the minimal additional information the
framework must obtain from an operator class for sorting-based construction to
be possible, and (ii)~what must the framework do with that information so that
the resulting tree is no worse, for queries, than one built by insertion?

The answer to (i) turns out to be small: one optional support function
returning a comparator over keys. The order need not be total in any meaningful
sense, and it need not be derivable from the query semantics; it must merely
preserve locality, and only the operator class can know which order does.

The answer to (ii) is the substance of this paper, because the obvious
procedure is wrong in a way that measurement of construction time cannot
reveal. Filling pages to capacity and starting a new page chooses split points
by a single criterion~--- occupancy. In one dimension that criterion is
optimal. In more than one it is not: an order that preserves locality does so
only approximately, and a page boundary that falls inside a dense region
produces two neighbouring pages whose keys overlap substantially. Overlap
determines the cost of search in a generalized tree: a query descends into
every subtree whose key matches it. Measured per level, the overlap of an index
built this way settles at about four~--- four subtrees entered at each level
instead of one~--- and stays there as the tree grows deeper, whereas
partitioning by the operator class settles at about $1.3$. The index is built
several times faster, occupies less space, and answers point queries twice as
expensively, consistently across an order of magnitude in data volume. A
comparison of construction methods by construction time~--- the comparison such
papers customarily report~--- hides this entirely.

\paragraph{Contributions.}

\begin{enumerate}
\item We formulate bulk loading as a framework-level problem and show that a
      single comparator-returning support function is sufficient, so that the
      technique becomes available to an arbitrary operator class rather than to
      one key type (Section~\ref{sec:build}). Changing the space-filling order
      requires writing one function and touches neither the tree, nor logging,
      nor concurrency: we demonstrate this by supplying Hilbert order as a
      drop-in alternative to Z-order, an operator class differing from the
      built-in one in exactly one support function out of ten
      (Section~\ref{sec:eval}).
\item We show that the order and the choice of split points are two separate
      levers, both of which live in the operator class, and we measure their
      interaction. On synthetic data they are partially interchangeable: a
      better curve recovers 44\,\% of the penalty of occupancy-based packing, a
      better split rule recovers 53\,\%, and together they bring query cost
      within 4\,\% of insertion-based construction. On 138M real points the
      relation is asymmetric: the curve recovers 33\,\% when split points are
      chosen by occupancy and nothing at all when they are chosen by the
      operator class. Investing in an elaborate space-filling order therefore
      pays off only while the partitioning rule stays naive
      (Section~\ref{sec:eval}).
\item We characterise the overlap introduced by occupancy-based packing. It does
      not compound with depth, as one might expect from the fact that a node is
      visited only if its parent was; instead each strategy converges to a
      steady-state value $\omega_\infty$ within one or two levels of the root.
      We measure $\omega_\infty$ for both strategies, across data volume and
      degree of clustering (Section~\ref{sec:overlap}).
\item We give an algorithm that delegates split point selection back to the
      operator class over a buffer of $k$ pages, removing 87--89\,\% of the
      penalty, and analyse the dependence on $k$ (Section~\ref{sec:overlap}).
      The residual is the price of choosing split points within a buffer rather
      than globally, and we quantify it against the insertion-built baseline.
\item We report five years of deployment: the design ships in PostgreSQL since
      release~14 and is the default for a family of operator classes since
      release~18. We include the one case in which the correctness condition of
      the method was violated and what it implies for framework design
      (Section~\ref{sec:deployment}).
\end{enumerate}

